\documentclass[../../main.tex]{subfiles}
\begin{document}
\section{零测集}\label{section:ch06-零测集}

回忆集合 $A\subseteq\mathbb{R}^n$ 具有零测度的含义(见附录 C):对每个 $\delta>0$, $A$ 可被可数多个开长方体覆盖,其体积之和小于 $\delta$(\cref{figure:ch06-measure-zero})。

\begin{figure}[ht]
\centering
\includegraphics[width=0.8\textwidth]{fig-06-1.png}
\caption{零测集}\label{figure:ch06-measure-zero}
\end{figure}

\begin{exercise}\label{exercise:ch06-rectangles-balls-cubes}
证明在零测集定义中,开长方体可换成开球或开立方体。[提示:对立方体,先证明每个开长方体 $R\subseteq\mathbb{R}^n$ 可被有限多个开立方体覆盖,其体积和不超过 $R$ 体积的 $2^n$ 倍。对球,(\cref{figure:ch06-balls-cubes})提示了主要思想。]
\end{exercise}

\begin{figure}[ht]
\centering
\includegraphics[width=0.8\textwidth]{fig-06-2.png}
\caption{球与立方体}\label{figure:ch06-balls-cubes}
\end{figure}

我们需要关于零测集的一个技术引理。若你熟悉 Lebesgue 测度理论,会注意到这个结果可由乘积集上可测函数积分的 Fubini 定理轻松推出;但下面给出一个不依赖测度论的初等证明。

\begin{lemma}\label{lemma:ch06-compact-hyperplane}
设 $A\subseteq\mathbb{R}^n$ 是紧子集,且对每个 $c\in\mathbb{R}$, $A$ 与 $\{c\}\times\mathbb{R}^{n-1}$ 的交具有 $(n-1)$ 维零测度。则 $A$ 具有 $n$ 维零测度。
\end{lemma}

\begin{proof}
选取区间 $[a,b]\subseteq\mathbb{R}$ 使 $A\subseteq[a,b]\times\mathbb{R}^{n-1}$。对每个 $c\in[a,b]$,令 $A_c\subseteq\mathbb{R}^{n-1}$ 表示紧子集
\begin{align*}
\left\{x\in\mathbb{R}^{n-1}:(c,x)\in A\right\}.
\end{align*}
给定 $\delta>0$。假设蕴涵:对每个 $c\in[a,b]$,集合 $A_c$ 可被有限多个 $(n-1)$ 维开立方体 $C_1,\dots,C_k$ 覆盖,其总体积小于 $\delta$。令 $U_c=C_1\cup\cdots\cup C_k\subseteq\mathbb{R}^{n-1}$。由于 $A$ 紧,必存在含 $c$ 的开区间 $J_c$,使 $A$ 与 $J_c\times\mathbb{R}^{n-1}$ 的交含于 $J_c\times U_c$;否则将存在一列点 $(c_i,x_i)\in A$,使 $c_i\to c$ 且 $x_i\notin U_c$;但取收敛子列后得到 $x_i\to x\in A_c\setminus U_c$,这与 $A_c\subseteq U_c$ 矛盾。区间族 $\{J_c:c\in[a,b]\}$ 构成 $[a,b]$ 的开覆盖,故存在有限多个数 $c_1<\cdots<c_m$,使区间 $J_{c_1},\dots,J_{c_m}$ 覆盖 $[a,b]$。必要时缩小重叠处的 $J_{c_i}$,可使 $J_{c_1},\dots,J_{c_m}$ 的长度之和不超过 $2|b-a|$。于是 $A$ 含于
\begin{align*}
(J_{c_1}\times U_{c_1})\cup\cdots\cup(J_{c_m}\times U_{c_m}),
\end{align*}
它是有限多个开长方体之并,其总体积小于 $2\delta|b-a|$。由于它可取得任意小,故 $A$ 具有 $n$ 维零测度。
\end{proof}

最重要的零测集是连续函数的图。

\begin{proposition}\label{proposition:ch06-graph-measure-zero}
设 $A$ 是 $\mathbb{R}^{n-1}$ 或 $\mathbb{H}^{n-1}$ 的开子集或闭子集, $f:A\to\mathbb{R}$ 是连续函数。则 $f$ 的图在 $\mathbb{R}^n$ 中零测。
\end{proposition}

\begin{proof}
先设 $A$ 紧。在此情形我们对 $n$ 归纳。当 $n=1$ 时, $f$ 的图至多是单点,结论平凡。为作归纳步,我们应用\cref{lemma:ch06-compact-hyperplane}。对每个 $c\in\mathbb{R}$, $f$ 的图与 $\{c\}\times\mathbb{R}^{n-1}$ 的交恰为 $f$ 在 $\{x\in A:x^1=c\}$ 上限制的图,而后者又是 $n-2$ 个变量的连续函数的图。由归纳假设,每个这样的图具有 $(n-1)$ 维零测度,故由\cref{lemma:ch06-compact-hyperplane}, $f$ 的图本身具有 $n$ 维零测度。若 $A$ 非紧,则由\cref{proposition:appA-sigma-compact} 知它是紧子集的可数并,故 $f$ 的图是零测集的可数并。
\end{proof}

\begin{corollary}\label{corollary:ch06-affine-subspace-measure-zero}
$\mathbb{R}^n$ 的每个真仿射子空间在 $\mathbb{R}^n$ 中零测。
\end{corollary}

\begin{proof}
设 $S\subseteq\mathbb{R}^n$ 是真仿射子空间。先设 $\dim S=n-1$。则至少有一条坐标轴,比如 $x^i$ 轴,不平行于 $S$;此时 $S$ 是形如
\begin{align*}
x^i=F\bigl(x^1,\dots,x^{i-1},x^{i+1},\dots,x^n\bigr)
\end{align*}
的仿射函数的图,故由\cref{proposition:ch06-graph-measure-zero} 知它零测。若 $\dim S<n-1$,则 $S$ 含于某 $n-1$ 维仿射子空间中,故由\cref{proposition:appC-subspace-measure}(b) 知 $S$ 零测。
\end{proof}

我们的目标是把零测的概念以微分同胚不变的方式推广到流形的子集上。由于流形不带度量,立方体或球的体积没有意义,故不能直接套用相同定义。然而,关键由下面的命题提供,它蕴涵零测这一性质对 $\mathbb{R}^n$ 的子集而言是微分同胚不变的。

\begin{proposition}\label{proposition:ch06-zero-measure-smooth-image}
设 $A\subseteq\mathbb{R}^n$ 零测, $F:A\to\mathbb{R}^n$ 是光滑映射。则 $F(A)$ 零测。
\end{proposition}

\begin{proof}
由定义,对每个 $p\in A$, $F$ 可延拓为 $\mathbb{R}^n$ 中 $p$ 的某邻域上的光滑映射,仍记为 $F$。必要时缩小该邻域,可设存在含 $p$ 的开球 $U$,使 $F$ 光滑延拓到 $\overline{U}$ 上。由\cref{proposition:appA-second-countable}, $A$ 被可数多个这样的预紧开子集覆盖,故 $F(A)$ 是形如 $F(A\cap\overline{U})$ 的可数多个集合之并。于是只需证明每个这样的集合零测。

由于 $\overline{U}$ 紧,存在常数 $C$,使对所有 $x\in\overline{U}$ 有 $|\mathrm{D}F(x)|\leqslant C$。利用光滑函数的 Lipschitz 估计(\cref{proposition:appC-lipschitz}),对所有 $x,x'\in\overline{U}$ 有
\begin{align}\label{align:ch06-lipschitz}
|F(x)-F(x')|\leqslant C|x-x'|.
\end{align}
给定 $\delta>0$,用开球 $\{B_j\}$ 可数覆盖 $A\cap\overline{U}$,使 $\sum_j\operatorname{Vol}(B_j)<\delta$。由\eqref{align:ch06-lipschitz}, $F(\overline{U}\cap B_j)$ 含于半径不超过 $C$ 倍 $B_j$ 半径的球 $\widetilde{B}_j$ 中(\cref{figure:ch06-image-measure-zero})。于是 $F(A\cap\overline{U})$ 含于球族 $\{\widetilde{B}_j\}$ 之并,其体积和至多为 $C^n\delta$。由于它可取得任意小,故 $F(A\cap\overline{U})$ 零测。
\end{proof}

\begin{figure}[ht]
\centering
\includegraphics[width=0.8\textwidth]{fig-06-3.png}
\caption{零测集的像}\label{figure:ch06-image-measure-zero}
\end{figure}

若 $M$ 是 $n$ 维带边或不带边光滑流形,我们说子集 $A\subseteq M$ 在 $M$ 中零测,如果对 $M$ 的每个光滑坐标卡 $(U,\varphi)$,子集 $\varphi(A\cap U)\subseteq\mathbb{R}^n$ 具有 $n$ 维零测度。下面的引理表明,只需对其定义域覆盖 $A$ 的某一组光滑坐标卡检验这一条件即可。

\begin{lemma}\label{lemma:ch06-chart-zero-measure}
设 $M$ 是 $n$ 维带边或不带边光滑流形, $A\subseteq M$。设存在一组光滑坐标卡 $\{(U_\alpha,\varphi_\alpha)\}$,其定义域覆盖 $A$,且对每个 $\alpha$, $\varphi_\alpha(A\cap U_\alpha)$ 在 $\mathbb{R}^n$ 中零测。则 $A$ 在 $M$ 中零测。
\end{lemma}

\begin{proof}
任取光滑坐标卡 $(V,\psi)$。我们需证明 $\psi(A\cap V)$ 零测。可数多个 $U_\alpha$ 覆盖 $A\cap V$。对每个这样的 $U_\alpha$,有
\begin{align*}
\psi(A\cap V\cap U_\alpha)=\bigl(\psi\circ\varphi_\alpha^{-1}\bigr)\bigl(\varphi_\alpha(A\cap V\cap U_\alpha)\bigr).
\end{align*}
(见\cref{figure:ch06-proof-zero-chart}。)现在, $\varphi_\alpha(A\cap V\cap U_\alpha)$ 是 $\varphi_\alpha(A\cap U_\alpha)$ 的子集,由假设它在 $\mathbb{R}^n$ 中零测。故由\cref{proposition:ch06-zero-measure-smooth-image}(应用于 $\psi\circ\varphi_\alpha^{-1}$), $\psi(A\cap V\cap U_\alpha)$ 零测。由于 $\psi(A\cap V)$ 是可数多个这样集合之并,它也零测。
\end{proof}

\begin{figure}[ht]
\centering
\includegraphics[width=0.8\textwidth]{fig-06-4.png}
\caption{引理 6.6 的证明}\label{figure:ch06-proof-zero-chart}
\end{figure}

\begin{exercise}\label{exercise:ch06-countable-union-zero}
设 $M$ 是带边或不带边光滑流形。证明 $M$ 中零测集的可数并仍零测。
\end{exercise}

\begin{proposition}\label{proposition:ch06-zero-complement-dense}
设 $M$ 是带边或不带边光滑流形, $A\subseteq M$ 在 $M$ 中零测。则 $M\setminus A$ 在 $M$ 中稠密。
\end{proposition}

\begin{proof}
若 $M\setminus A$ 不稠密,则 $A$ 含 $M$ 的某非空开子集;这意味着存在光滑坐标卡 $(V,\psi)$,使 $\psi(A\cap V)$ 含 $\mathbb{R}^n$(其中 $n=\dim M$)的某非空开子集。但 $\psi(A\cap V)$ 在 $\mathbb{R}^n$ 中零测,这与\cref{corollary:appC-zero-open} 矛盾。
\end{proof}

\begin{theorem}\label{theorem:ch06-zero-measure-smooth-map}
设 $M$ 与 $N$ 是 $n$ 维带边或不带边光滑流形, $F:M\to N$ 是光滑映射, $A\subseteq M$ 是零测子集。则 $F(A)$ 在 $N$ 中零测。
\end{theorem}

\begin{proof}
设 $\{(U_i,\varphi_i)\}$ 是 $M$ 的可数光滑坐标覆盖。我们需证明:对 $N$ 的每个光滑坐标卡 $(V,\psi)$,集合 $\psi(F(A)\cap V)$ 在 $\mathbb{R}^n$ 中零测。注意该集合是可数多个形如
\begin{align*}
\psi\circ F\circ\varphi_i^{-1}\bigl(\varphi_i(A\cap U_i\cap F^{-1}(V))\bigr)
\end{align*}
的子集之并,由\cref{proposition:ch06-zero-measure-smooth-image} 知其中每个都零测。
\end{proof}

\end{document}